\chapter{Thực nghiệm và Đánh giá}

Chương này trình bày kết quả thực nghiệm của hệ thống đề xuất trên 14 ca
kiểm thử, gồm tám ứng dụng lab và sáu CVE thực tế.
Phần~\ref{sec:exp-env} mô tả môi trường và thiết lập thực nghiệm.
Phần~\ref{sec:exp-targets} liệt kê toàn bộ mục tiêu.
Phần~\ref{sec:exp-lab} trình bày kết quả nhóm lab.
Phần~\ref{sec:exp-cve} phân tích từng CVE.
Phần~\ref{sec:exp-summary} tổng hợp và đánh giá.

% ============================================================
\section{Môi trường Thực nghiệm}
\label{sec:exp-env}
% ============================================================

\subsection{Cấu hình phần cứng và phần mềm}

Toàn bộ thực nghiệm được thực hiện trên một máy tính với hệ điều hành
Ubuntu 22.04 LTS, CPU Intel Core i7 thế hệ 12, RAM 32 GB. Mỗi ca kiểm
thử được khởi tạo độc lập bằng lệnh:

\begin{verbatim}
docker compose -f docker-compose.<target>.yml up --build
\end{verbatim}

đảm bảo trạng thái sạch (clean state) trước mỗi lần chạy. Giới hạn
thời gian tối đa là 600 giây mỗi ca kiểm thử. Bảng~\ref{tab:env} tóm
tắt cấu hình môi trường.

\begin{table}[h]{Cấu hình môi trường thực nghiệm}
\label{tab:env}
\centering
\begin{tabular}{|l|l|}
\hline
\textbf{Thành phần} & \textbf{Chi tiết} \\
\hline
Hệ điều hành    & Ubuntu 22.04 LTS \\
CPU             & Intel Core i7 thế hệ 12 \\
RAM             & 32 GB \\
Docker          & Docker Engine 24.x \\
PHP web server  & Apache 2.4, PHP 7.4 / 8.2 (tùy target) \\
Database        & MySQL 8 (khi cần) \\
Giới hạn thời gian & 600 giây/ca kiểm thử \\
Subnet Docker   & \texttt{172.20.0.0/24}, OOB server IP \texttt{172.20.0.10} \\
\hline
\end{tabular}
\end{table}

\subsection{Cấu hình Fuzzer}

Hệ thống sử dụng Phuzz mở rộng với \texttt{SSRFMutator} và
\texttt{SSRFVulnCheck} đã tích hợp như mô tả ở Chương~3. Hai chế độ
hoạt động được dùng tùy loại target:

\begin{itemize}
    \item \texttt{"ssrf\_only": true} --- toàn bộ đột biến dùng
    \texttt{SSRFMutator} (21 biến thể OOB). Áp dụng cho các endpoint
    nhận URL hoặc đường dẫn trực tiếp.
    \item \texttt{"cmdi\_ssrf\_only": true} --- toàn bộ đột biến dùng
    \texttt{CmdInjectionSSRFMutator} (10 biến thể shell injection). Áp
    dụng cho các endpoint truyền tham số vào lệnh shell.
\end{itemize}

21 biến thể của \texttt{SSRFMutator} được tổ chức thành ba nhóm: biến
thể IP-based (v0--v14, sử dụng IP 172.20.0.10 dưới các dạng mã hóa
khác nhau), biến thể DNS-based (v15--v18, token nằm trong subdomain và
được phân giải qua dnsmasq wildcard \texttt{*.oob}), và biến thể
hostname-based (v19--v20, dùng hostname \texttt{oob.phuzz} được phân
giải qua \texttt{extra\_hosts} trong Docker).

\subsection{Đường cơ sở so sánh}

Mỗi ca kiểm thử được chạy với hai cấu hình để xác định đóng góp của
hệ thống đề xuất:

\begin{itemize}
    \item \textbf{Phuzz gốc (P)}: không có \texttt{SSRFMutator} /
    \texttt{SSRFVulnCheck}; chỉ dùng tín hiệu coverage và exception.
    \item \textbf{Phuzz+SSRF (P+S)}: bản mở rộng đầy đủ với tín hiệu
    OOB.
\end{itemize}

Kết quả của Phuzz gốc trên toàn bộ 14 ca: \textbf{0 lỗ hổng SSRF được
phát hiện} --- xác nhận thực nghiệm rằng tín hiệu exception và coverage
không đủ để phát hiện Blind SSRF.

% ============================================================
\section{Danh sách Mục tiêu Thực nghiệm}
\label{sec:exp-targets}
% ============================================================

Bảng~\ref{tab:lab-targets} và Bảng~\ref{tab:cve-targets} liệt kê toàn
bộ 14 mục tiêu, phân loại theo nhóm lab và CVE.

\begin{table}[h]{Danh sách ứng dụng lab được kiểm thử}
\label{tab:lab-targets}
\centering
\small
\begin{tabular}{|c|l|l|l|l|}
\hline
\textbf{STT} & \textbf{Ứng dụng} & \textbf{Endpoint (param)} & \textbf{Sink} & \textbf{Mutator} \\
\hline
L1 & XVWA \texttt{ssrf\_xspa}
   & \texttt{POST img\_url=}
   & \texttt{curl\_exec}
   & SSRFMutator \\
L2 & Pikachu SSRF (2 ep.)
   & \texttt{GET ?url=} / \texttt{GET ?file=}
   & \texttt{curl\_exec} / \texttt{fgc}
   & SSRFMutator \\
L3 & ssrf-lab (2 ep.)
   & \texttt{POST handler=} (basic + bypass)
   & \texttt{curl\_exec}
   & SSRFMutator \\
L4 & SSRF\_Vul\_Lab (2 ep.)
   & \texttt{POST file=} (fgc + dns\_rebinding)
   & \texttt{file\_get\_contents}
   & SSRFMutator \\
L5 & BTSLab SSRF
   & \texttt{POST url=}
   & \texttt{file\_get\_contents}
   & SSRFMutator \\
L6 & DVWA Cmd Inject.
   & \texttt{POST ip=}
   & \texttt{shell\_exec}
   & CmdInjSSRF \\
L7 & bWAPP RFI
   & \texttt{GET ?language=}
   & \texttt{include(http://)}
   & SSRFMutator \\
L8 & Mutillidae RFI
   & \texttt{GET ?page=}
   & \texttt{require\_once(http://)}
   & SSRFMutator \\
\hline
\end{tabular}
\end{table}

\begin{table}[h]{Sáu CVE được kiểm thử}
\label{tab:cve-targets}
\centering
\small
\begin{tabular}{|c|l|l|l|l|}
\hline
\textbf{STT} & \textbf{CVE} & \textbf{Ứng dụng} & \textbf{Endpoint} & \textbf{Sink} \\
\hline
C1 & CVE-2020-7071
   & PHP \texttt{filter\_var} app
   & \texttt{GET ?url=}
   & \texttt{file\_get\_contents} \\
C2 & CVE-2020-24148
   & Import XML Feed (WP)
   & \texttt{POST data=} (AJAX)
   & \texttt{file\_get\_contents} \\
C3 & CVE-2020-28975
   & Canto 1.3.0 \texttt{tree.php}
   & \texttt{GET ?subdomain=}
   & \texttt{curl\_exec} \\
C4 & CVE-2021-32682
   & elFinder $<$ 2.1.59
   & \texttt{GET ?url=}
   & \texttt{file\_get\_contents} \\
C5 & CVE-2022-1751
   & WP Skitter Slideshow
   & \texttt{GET ?image=}
   & \texttt{getimagesize} \\
C6 & CVE-2023-2249
   & wpForo (WP)
   & \texttt{POST image\_blob=}
   & \texttt{file\_get\_contents} \\
\hline
\end{tabular}
\end{table}

Lưu ý: CVE-2020-28975 thuộc plugin Canto 1.3.0 có ba file PHP đều dễ
bị tấn công theo cùng cơ chế --- \texttt{tree.php} (CVE-2020-28975),
\texttt{detail.php} (CVE-2020-28976), \texttt{get.php}
(CVE-2020-28977). Thực nghiệm kiểm thử đại diện qua \texttt{tree.php};
hai file còn lại chia sẻ cùng vector tấn công và không cần kiểm thử
riêng. CVE-2021-21311 (Adminer Elasticsearch) không được đưa vào do
Adminer throw \texttt{TypeError} trên PHP 8.2 trước khi đạt tới code
path SSRF.

% ============================================================
\section{Kết quả Nhóm Lab}
\label{sec:exp-lab}
% ============================================================

\subsection{Kết quả tổng hợp}

Bảng~\ref{tab:lab-results} trình bày kết quả phát hiện cho tám ứng dụng
lab, bao gồm 11 endpoint-level test run (ba ứng dụng L2, L3, L4 mỗi
ứng dụng có hai endpoint riêng biệt).

\begin{table}[h]{Kết quả phát hiện SSRF trên nhóm lab}
\label{tab:lab-results}
\centering
\small
\begin{tabular}{|c|l|l|l|c|c|}
\hline
\textbf{ID} & \textbf{Ứng dụng} & \textbf{Biến thể phát hiện} & \textbf{Hook function} & \textbf{P+S} & \textbf{P} \\
\hline
L1 & XVWA ssrf\_xspa
   & v0 (plain IP)
   & \texttt{curl\_exec}
   & \checkmark & \ding{55} \\
L2 & Pikachu (2 ep.)
   & v0 (cả 2 endpoint)
   & \texttt{curl\_exec} / \texttt{fgc}
   & \checkmark & \ding{55} \\
L3 & ssrf-lab (2 ep.)
   & v0 (basic), v1 (bypass)
   & \texttt{curl\_exec}
   & \checkmark & \ding{55} \\
L4 & SSRF\_Vul\_Lab (2 ep.)
   & v0 (fgc), v15 (dns)
   & \texttt{file\_get\_contents}
   & \checkmark & \ding{55} \\
L5 & BTSLab SSRF
   & v0 (plain IP)
   & \texttt{file\_get\_contents}
   & \checkmark & \ding{55} \\
L6 & DVWA Cmd Inject.
   & Ck (semicolon)
   & \texttt{shell\_exec}
   & \checkmark & \ding{55} \\
L7 & bWAPP RFI
   & v0 (plain IP)
   & stream wrapper \texttt{http://}
   & \checkmark & \ding{55} \\
L8 & Mutillidae RFI
   & v0 (plain IP)
   & stream wrapper \texttt{http://}
   & \checkmark & \ding{55} \\
\hline
\multicolumn{4}{|l|}{\textbf{Tổng (8 ứng dụng / 11 endpoint)}}
   & \textbf{8/8} & \textbf{0/8} \\
\hline
\end{tabular}
\end{table}

\subsection{Phân tích từng ứng dụng}

\textbf{XVWA (L1).} XVWA SSRF XSPA (Cross-Site Port Attack) là mô-đun
SSRF của ứng dụng eXtremely Vulnerable Web Application, trong đó tham
số \texttt{img\_url} được truyền trực tiếp vào \texttt{curl\_exec} mà
không có bất kỳ kiểm tra nào. Biến thể v0 (plain IP) kích hoạt hook
ngay trong vòng lặp đầu tiên của fuzzer.

\textbf{Pikachu (L2).} Pikachu có hai module SSRF riêng biệt đóng vai
trò sanity check: \texttt{ssrf\_curl.php} nhận tham số \texttt{url} rồi
truyền vào \texttt{curl\_exec}, và \texttt{ssrf\_fgc.php} nhận tham số
\texttt{file} rồi truyền vào \texttt{file\_get\_contents}. Cả hai
endpoint không có cơ chế lọc, đều phát hiện với v0. Trường hợp này xác
nhận rằng hai sink function phổ biến nhất (\texttt{curl\_exec} và
\texttt{file\_get\_contents}) đều được hook và phát hiện đúng.

\textbf{ssrf-lab (L3).} ssrf-lab cung cấp hai endpoint trên cùng file
\texttt{testhook.php}: endpoint \textit{basic} (không có filter) và
endpoint \textit{advanced1} (có filter đơn giản dựa trên chuỗi IP).
Endpoint basic phát hiện với v0. Endpoint advanced1 chặn các chuỗi chứa
dạng IP thập phân phân cách dấu chấm (\texttt{172.x.x.x}); biến thể v1
(decimal encoding: \texttt{http://2886795274:8000/<token>}, trong đó
2886795274 = 172.20.0.10 dưới dạng số nguyên 32-bit) vượt qua filter và
kích hoạt \texttt{curl\_exec}. Trường hợp này minh họa cơ chế bypass
cơ bản nhất: filter kiểm tra chuỗi thay vì kiểm tra giá trị IP đã chuẩn
hóa.

\textbf{SSRF\_Vulnerable\_Lab (L4).} Ứng dụng có hai endpoint:
\texttt{file\_get\_content.php} (endpoint fgc, không có filter) và
\texttt{dns\_rebinding.php} (mô phỏng tình huống có cơ chế kiểm tra
hostname). Endpoint fgc phát hiện với v0. Endpoint dns\_rebinding sử
dụng biến thể v15 --- payload dạng \texttt{http://<token>.oob:8000/},
trong đó token nằm trong subdomain; dnsmasq trong Docker network phân
giải wildcard \texttt{*.oob $\rightarrow$ 172.20.0.10}, cho phép OOB
callback đến đích mà không cần truyền IP nội bộ trực tiếp trong URL.
Endpoint này được cấu hình với \texttt{request\_timeout: 15} giây để
chờ DNS resolution hoàn tất.

\textbf{BTSLab (L5).} BTSLab yêu cầu đăng nhập trước khi truy cập
endpoint SSRF. Fuzzer sử dụng module \texttt{btslab\_requests} để tự
động đăng nhập và lưu session cookie \texttt{PHPSESSID}, sau đó đột
biến tham số \texttt{url} tại \texttt{/vulnerability/ssrf/ssrf.php}.
Sink function là \texttt{file\_get\_contents}; phát hiện với v0.

\textbf{DVWA Command Injection (L6).} Đây là trường hợp SSRF gián tiếp:
endpoint thực thi \texttt{shell\_exec('ping -c 4 ' . \$target)} ---
tham số \texttt{ip} do người dùng kiểm soát được nhúng trực tiếp vào
lệnh shell. \texttt{CmdInjectionSSRFMutator} sinh payload semicolon:
\texttt{127.0.0.1; curl http://172.20.0.10:8000/<token>}, khiến shell
thực thi lệnh \texttt{curl} sau \texttt{ping}. Hook trên
\texttt{shell\_exec} trong \texttt{08\_ssrf.php} ghi nhận lời gọi; OOB
server nhận callback từ \texttt{curl}. Trường hợp này mở rộng phạm vi
phát hiện sang lớp lỗ hổng kết hợp Command Injection và SSRF.

\textbf{bWAPP RFI (L7) và Mutillidae RFI (L8).} Hai trường hợp Remote
File Inclusion sử dụng language construct của PHP: \texttt{include(\$language)}
trong bWAPP (\texttt{rlfi.php}, security level 0) và
\texttt{require\_once(\$page)} trong Mutillidae (\texttt{index.php},
security level 0). Language construct là cấu trúc ngôn ngữ, không thể
hook bằng \texttt{uopz\_set\_hook}. Thay vào đó,
\texttt{09\_ssrf\_stream.php} đăng ký lại PHP stream wrapper
\texttt{http://}, chặn mọi lần PHP mở URL qua giao thức HTTP, kể cả từ
\texttt{include} và \texttt{require\_once}. Biến thể v0 kích hoạt stream
wrapper và OOB callback nhận được. Kết quả xác nhận lớp phát hiện qua
stream wrapper là thành phần bổ sung không thể thiếu bên cạnh hook UOPZ.

% ============================================================
\section{Kết quả Nhóm CVE}
\label{sec:exp-cve}
% ============================================================

\subsection{Kết quả tổng hợp}

Bảng~\ref{tab:cve-results} tóm tắt kết quả phát hiện cho sáu CVE.

\begin{table}[h]{Kết quả phát hiện SSRF trên nhóm CVE}
\label{tab:cve-results}
\centering
\small
\begin{tabular}{|c|l|l|l|c|c|}
\hline
\textbf{ID} & \textbf{CVE} & \textbf{Biến thể phát hiện} & \textbf{Auth} & \textbf{P+S} & \textbf{P} \\
\hline
C1 & CVE-2020-7071  & v0 (plain IP)        & Không  & \checkmark & \ding{55} \\
C2 & CVE-2020-24148 & v0 (plain IP)        & Admin  & \checkmark & \ding{55} \\
C3 & CVE-2020-28975 & v20 (scheme-less)    & Không  & \checkmark & \ding{55} \\
C4 & CVE-2021-32682 & v19 (hostname)       & Không  & \checkmark & \ding{55} \\
C5 & CVE-2022-1751  & v0 (plain IP)        & Không  & \checkmark & \ding{55} \\
C6 & CVE-2023-2249  & v0 (plain IP)        & User   & \checkmark & \ding{55} \\
\hline
\multicolumn{4}{|l|}{\textbf{Tổng}} & \textbf{6/6} & \textbf{0/6} \\
\hline
\end{tabular}
\end{table}

\subsection{CVE-2020-7071 --- Bypass \texttt{filter\_var} trong PHP}

Ứng dụng kiểm tra URL đầu vào bằng \texttt{filter\_var(\$url,
FILTER\_VALIDATE\_URL)} trước khi truyền vào
\texttt{file\_get\_contents}. Trên PHP 7.4.13 (phiên bản có lỗi
CVE-2020-7071), \texttt{filter\_var} chấp nhận URL dạng
\texttt{http://172.20.0.10/} là hợp lệ --- đây chính là nội dung của
lỗ hổng: hàm kiểm tra URL không phân biệt địa chỉ nội bộ và địa chỉ
công khai. Biến thể v0 (plain IP) vượt qua được kiểm tra này và kích
hoạt OOB callback. Container web sử dụng \texttt{Dockerfile.php74} để
đảm bảo đúng phiên bản PHP 7.4 có lỗi.

\subsection{CVE-2020-24148 --- Import XML Feed (WordPress Plugin)}

Plugin Import XML Feed $\leq$ 2.0.0 có endpoint AJAX
\texttt{moove\_read\_xml} yêu cầu quyền admin nhưng không kiểm tra
nonce. Fuzzer đăng nhập với tài khoản admin thông qua module session
management của Phuzz, rồi đột biến tham số \texttt{data} với payload
OOB. Sink function là \texttt{file\_get\_contents}; phát hiện với v0.
Thời gian phát hiện đạt \textbf{$\approx$506 giây} --- cao nhất trong
toàn bộ thực nghiệm, chủ yếu do overhead từ WordPress bootstrap (khởi
động MySQL và WordPress init sequence) và quá trình xác thực admin. Với
giới hạn 600 giây, biên độ an toàn không lớn; đây cũng là một trong
các hạn chế được thảo luận ở Chương~5.

\subsection{CVE-2020-28975 --- Canto 1.3.0 WordPress Plugin}

Plugin Canto 1.3.0 xây dựng URL theo mẫu cứng
\texttt{https://\{subdomain\}.\{app\_api\}/api/v1/...} rồi gọi cURL
với \texttt{CURLOPT\_SSL\_VERIFYPEER = false}. Đây là trường hợp đặc
biệt: các biến thể IP-based (v0--v14) đều không hoạt động vì IP được
nhúng vào vị trí subdomain tạo ra URL dạng
\texttt{https://172.20.0.10.x/...} --- chuỗi này không hợp lệ về mặt
DNS và cURL thất bại ở bước phân giải.

Hệ thống phát hiện qua \textbf{biến thể v20} (scheme-less:
\texttt{oob.phuzz:8443/<token>\#}): khi Canto nhận
\texttt{subdomain = "oob.phuzz:8443/<token>\#"}, URL được xây dựng
thành \texttt{https://oob.phuzz:8443/<token>\#.x/api/v1/...}. cURL
phân tích \texttt{oob.phuzz:8443} là host và \texttt{/<token>} là path;
phần sau \texttt{\#} là fragment bị bỏ qua. Entry \texttt{extra\_hosts}
trong Docker phân giải \texttt{oob.phuzz $\rightarrow$ 172.20.0.10}; OOB
HTTPS listener (cổng 8443, SSL verification disabled) nhận callback và
ghi log.

\subsection{CVE-2021-32682 --- elFinder $<$ 2.1.59}

elFinder triển khai kiểm tra địa chỉ IP nội bộ (RFC~1918) trong hàm
\texttt{get\_remote\_contents()} bằng regex trên chuỗi URL. Các biến
thể IP trực tiếp (plain, decimal, hex, octal, IPv6) đều bị chặn tại
tầng kiểm tra chuỗi này trước khi đến \texttt{file\_get\_contents}.

Hệ thống phát hiện qua \textbf{biến thể v19}
(\texttt{http://oob.phuzz:8000/<token>}): hostname \texttt{oob.phuzz}
không khớp với bất kỳ pattern RFC~1918 nào trong regex của elFinder,
nên \texttt{file\_get\_contents} được gọi với URL này. Entry
\texttt{extra\_hosts} trong Docker phân giải
\texttt{oob.phuzz $\rightarrow$ 172.20.0.10} ở tầng OS; OOB server nhận
callback. Đây là ví dụ điển hình của kỹ thuật bypass
\textit{hostname indirection}: cơ chế lọc kiểm tra biểu diễn chuỗi của
URL thay vì kiểm tra địa chỉ IP thực sự sau DNS resolution.

\subsection{CVE-2022-1751 --- WP Skitter Slideshow}

Plugin \texttt{wp-skitter-slideshow} gọi
\texttt{getimagesize(\$\_GET['image'])} trực tiếp trên URL người dùng
cung cấp trong file \texttt{image.php} mà không qua WordPress bootstrap.
\texttt{getimagesize} là hàm xử lý ảnh, không phải hàm mạng điển hình,
nhưng PHP triển khai nó với khả năng fetch URL từ xa qua HTTP wrapper
khi đối số là URL bắt đầu bằng \texttt{http://}. Trường hợp này xác
nhận tầm quan trọng của việc hook đủ rộng: nhóm Network/HTTP trong
\texttt{08\_ssrf.php} bao gồm \texttt{getimagesize} bên cạnh
\texttt{file\_get\_contents} và các hàm mạng thông thường. Biến thể v0
kích hoạt OOB callback mà không yêu cầu xác thực. Plugin đã bị đóng
trên WordPress.org; mã nguồn lấy từ GitHub mirror.

\subsection{CVE-2023-2249 --- wpForo Profile Cover Upload}

Plugin wpForo $\leq$ 2.1.7 cho phép thành viên đã đăng nhập ở mức
Subscriber cập nhật ảnh bìa profile bằng URL. Endpoint AJAX
\texttt{wpforo\_upload\_remote\_image} fetch URL ảnh phía server thông
qua \texttt{file\_get\_contents} mà không kiểm tra địa chỉ đích. Fuzzer
đăng nhập với tài khoản subscriber (\texttt{subscriber1/subscriber1}),
rồi đột biến tham số \texttt{image\_blob} với v0. Đây là SSRF
\textit{authenticated} ở mức quyền thấp nhất, xác nhận hệ thống xử lý
đúng session management cấp user thông thường bên cạnh cấp admin.

% ============================================================
\section{Tổng hợp và Đánh giá}
\label{sec:exp-summary}
% ============================================================

\subsection{Kết quả tổng thể}

Bảng~\ref{tab:overall} tổng hợp kết quả trên toàn bộ 14 ca kiểm thử.

\begin{table}[h]{Tổng hợp kết quả theo nhóm}
\label{tab:overall}
\centering
\begin{tabular}{|l|c|c|c|c|}
\hline
\textbf{Nhóm} & \textbf{Tổng ca} & \textbf{P+S phát hiện} & \textbf{P phát hiện} & \textbf{FP} \\
\hline
Lab (8 ứng dụng, 11 endpoint) & 8  & 8 (100\%) & 0 & 0 \\
CVE                            & 6  & 6 (100\%) & 0 & 0 \\
\hline
\textbf{Tổng}                 & \textbf{14} & \textbf{14 (100\%)} & \textbf{0} & \textbf{0} \\
\hline
\end{tabular}
\end{table}

\begin{table}[h]{Phân bố sink function được kích hoạt}
\label{tab:sinks}
\centering
\begin{tabular}{|l|c|l|}
\hline
\textbf{Sink function} & \textbf{Số run} & \textbf{Nhóm hook} \\
\hline
\texttt{file\_get\_contents}                      & 8  & Network/HTTP (\texttt{08\_ssrf.php}) \\
\texttt{curl\_exec}                                & 5  & cURL (\texttt{08\_ssrf.php}) \\
\texttt{include} / \texttt{require\_once} (stream) & 2  & Stream wrapper (\texttt{09\_ssrf\_stream.php}) \\
\texttt{shell\_exec}                               & 1  & Shell Execution (\texttt{08\_ssrf.php}) \\
\texttt{getimagesize}                              & 1  & Network/HTTP (\texttt{08\_ssrf.php}) \\
\hline
\textbf{Tổng}                                     & \textbf{17} & \\
\hline
\end{tabular}
\end{table}

\begin{table}[h]{Phân bố biến thể bypass được sử dụng}
\label{tab:variants}
\centering
\begin{tabular}{|l|c|l|}
\hline
\textbf{Biến thể} & \textbf{Số run} & \textbf{Áp dụng cho} \\
\hline
v0  (plain IP \texttt{172.20.0.10})
  & 12 & L1, L2$\times$2, L3-basic, L4-fgc, L5, L7, L8, C1, C2, C5, C6 \\
v1  (decimal IP \texttt{2886795274})
  & 1  & L3-advanced1 (bypass filter chuỗi IP) \\
v15 (DNS: \texttt{<token>.oob:8000/})
  & 1  & L4-dns\_rebinding \\
v19 (hostname: \texttt{oob.phuzz:8000/<token>})
  & 1  & C4 (elFinder, bypass RFC~1918 regex) \\
v20 (scheme-less: \texttt{oob.phuzz:8443/<token>\#})
  & 1  & C3 (Canto, bypass URL template) \\
Ck  (semicolon cmd inj.)
  & 1  & L6 (DVWA \texttt{shell\_exec}) \\
\hline
\textbf{Tổng} & \textbf{17} & \\
\hline
\end{tabular}
\end{table}

\subsection{Phân tích kết quả}

\textbf{Tỷ lệ phát hiện và False Positive.}
Hệ thống đề xuất đạt tỷ lệ phát hiện \textbf{14/14 (100\%)} trên toàn
bộ 14 ca kiểm thử, với tỷ lệ false positive bằng 0. Mỗi báo cáo đều
có đủ hai file bằng chứng (\texttt{ssrf-hooks/<ID>.json} và
\texttt{oob-logs/<token>.json}), xác nhận ứng dụng đã thực sự thực
hiện kết nối mạng tới OOB server. Ngược lại, Phuzz gốc không phát hiện
được bất kỳ trường hợp nào (0/14), khẳng định thực nghiệm rằng tín
hiệu coverage và exception hoàn toàn bất lực trước Blind SSRF.

\textbf{Tầm quan trọng của đa biến thể bypass.}
Bảng~\ref{tab:variants} cho thấy biến thể plain IP (v0) chiếm đa số
(12/17 run), nhưng năm run còn lại chỉ phát hiện được qua biến thể
đặc thù: Canto (C3) cần v20 do cơ chế xây dựng URL template; elFinder
(C4) cần v19 để bypass kiểm tra chuỗi RFC~1918; ssrf-vul-lab
dns\_rebinding cần v15 (DNS-based) để hoạt động trong môi trường có
DNS resolution; ssrf-lab advanced1 cần v1 để bypass filter chuỗi IP;
DVWA cần \texttt{CmdInjectionSSRFMutator}. Nếu hệ thống chỉ có biến
thể plain IP, tỷ lệ phát hiện sẽ giảm từ 100\% xuống khoảng 71\%
(12/17 run) ở mức endpoint và còn thấp hơn ở mức ứng dụng.

\textbf{Tầm quan trọng của hook rộng.}
Bảng~\ref{tab:sinks} cho thấy năm loại sink function khác nhau được
kích hoạt. Đặc biệt, \texttt{getimagesize} (C5) và cặp
\texttt{include}/\texttt{require\_once} qua stream wrapper (L7, L8) là
những sink không trực quan mà hầu hết công cụ bảo mật không theo dõi.
Phát hiện CVE-2022-1751 là bằng chứng rõ ràng: nếu thiếu hook
\texttt{getimagesize}, lỗ hổng này sẽ bị bỏ sót hoàn toàn. Tương tự,
nếu chỉ có hook UOPZ mà thiếu stream wrapper, bWAPP (L7) và Mutillidae
(L8) --- hai trường hợp RFI qua language construct --- sẽ không được
phát hiện.

\textbf{SSRF gián tiếp qua Command Injection (L6).}
DVWA là ví dụ về SSRF gián tiếp: tham số đầu vào không đi vào hàm
mạng mà đi vào shell command. \texttt{CmdInjectionSSRFMutator} và hook
nhóm Shell Execution xử lý đúng kịch bản này, mở rộng phạm vi phát
hiện sang lớp lỗ hổng kết hợp Command Injection và SSRF.

\textbf{Xác thực đa mức.}
Thực nghiệm bao phủ ba mức xác thực: không cần đăng nhập (L1--L4, C1,
C3--C5), đăng nhập user thông thường (L5--L8, C6), và đăng nhập admin
(C2). Trong tất cả trường hợp yêu cầu xác thực, module session
management của Phuzz xử lý đúng quy trình đăng nhập và duy trì session
cho các request tiếp theo, không làm gián đoạn vòng lặp fuzzing.

\section{Tổng kết Chương}

Thực nghiệm trên 14 ca kiểm thử (8 ứng dụng lab với tổng cộng 11
endpoint, 6 CVE) cho thấy hệ thống đề xuất đạt tỷ lệ phát hiện 100\%
với 0 false positive, trong khi Phuzz gốc không phát hiện được bất kỳ
trường hợp nào. Kết quả đặc biệt đáng chú ý gồm: (1) Canto (C3) và
elFinder (C4) chỉ phát hiện được qua biến thể đặc thù (scheme-less và
hostname indirection), đòi hỏi tập biến thể đa dạng hơn plain IP; (2)
CVE-2022-1751 kích hoạt sink \texttt{getimagesize} --- hàm ảnh không
được hầu hết công cụ coi là điểm SSRF; (3) bWAPP (L7) và Mutillidae
(L8) chứng minh sự cần thiết của lớp phát hiện qua stream wrapper bên
cạnh hook UOPZ. Chương~5 tổng kết đóng góp và thảo luận hướng phát
triển.
